\chapter{重积分}

\newpage

\section{二重积分}

\newpage

\subsection{化为累次积分}

设 $f(x,y)$ 在 $xOy$ 平面有界闭区域 $D$ 上有定义, 且下面出现的积分都存在, 则:

\begin{enumerate}[label=(\arabic*)]
    \item 当
    \[
        D_1=\{(x,y)\mid a\leq x\leq b,\ y_1(x)\leq y\leq y_2(x)\}
    \]
    (称为 $x$ 型区域)时,
    \[
        \iint\limits_{D_1}f(x,y)\,\mathrm{d}x\,\mathrm{d}y
        =\int_a^b\mathrm{d}x\int_{y_1(x)}^{y_2(x)}f(x,y)\,\mathrm{d}y.
    \]

    \item 当
    \[
        D_2=\{(x,y)\mid c\leq y\leq d,\ x_1(y)\leq x\leq x_2(y)\}
    \]
    (称为 $y$ 型区域)时,
    \[
        \iint\limits_{D_2}f(x,y)\,\mathrm{d}x\,\mathrm{d}y
        =\int_c^d\mathrm{d}y\int_{x_1(y)}^{x_2(y)}f(x,y)\,\mathrm{d}x.
    \]
\end{enumerate}

因此, 对于一般区域, 如区域可分成若干个不相交的 $x$ 型区域和 $y$ 型区域的并, 则可分别计算这些区域上积分的值再相加. 进一步, 也为后续的换元提供了方向——化为 $x$ 型区域(或 $y$ 型区域). 特别地, 若 $D=D_1=D_2$, 则上述两式均成立, 即两累次积分相等, 但具体计算时用哪个, 还得由被积函数决定.

\begin{example}
设 $f(x,y)$ 在 $\mathbb{R}^2$ 上连续, $a>0$, 试交换下列累次积分的次序:
\begin{enumerate}[label=(\arabic*)]
    \item $\displaystyle\int_0^a\mathrm{d}y\int_{\frac{y}{2}}^{2y}f(x,y)\,\mathrm{d}x$;
    \item $\displaystyle\int_{-1}^1\mathrm{d}x\int_{x^2+x}^{x+1}f(x,y)\,\mathrm{d}y$;
    \item $\displaystyle\int_0^a\mathrm{d}y\int_{-y}^{\sqrt{\sqrt{y}}}f(x,y)\,\mathrm{d}x$;
    \item $\displaystyle\int_0^{2a}\mathrm{d}x\int_{\sqrt{2ax-x^2}}^{4ax}f(x,y)\,\mathrm{d}y$.
\end{enumerate}
\end{example}
\exampleblank

\begin{example}[\markstar]
设 $f(x,y)$ 在 $\mathbb{R}^2$ 上连续, $a>0$, 试交换下列累次积分的次序:
\begin{enumerate}[label=(\arabic*)]
    \item $\displaystyle\int_0^{\frac{\pi}{2}}\mathrm{d}\theta\int_0^{2a\cos\theta}f(r,\theta)\,\mathrm{d}r$;
    \item $\displaystyle\int_0^{\frac{\pi}{4}}\mathrm{d}\theta\int_0^{a\sec\theta}f(r,\theta)\,\mathrm{d}r$;
    \item $\displaystyle\int_{-\frac{\pi}{2}}^{\frac{\pi}{2}}\mathrm{d}\theta\int_0^{2a\cos\theta}f(r\cos\theta,r\sin\theta)r\,\mathrm{d}r$.
\end{enumerate}
\end{example}
\exampleblank

\begin{example}
证明下列等式:
\begin{enumerate}[label=(\arabic*)]
    \item
    \[
        \int_0^e\mathrm{d}y\int_1^{y^2}\dfrac{\ln x}{e^x}\,\mathrm{d}x
        +\int_e^{e^2}\mathrm{d}y\int_{\ln y}^{2}\dfrac{\ln x}{e^x}\,\mathrm{d}x
        =\int_1^2\ln x\,\mathrm{d}x;
    \]
    \item
    \[
        \int_a^b\mathrm{d}y\int_a^y(y-x)^n f(x)\,\mathrm{d}x
        =\dfrac{1}{n+1}\int_a^b(b-x)^{n+1}f(x)\,\mathrm{d}x.
    \]
\end{enumerate}
\end{example}
\exampleblank

\begin{theorem}[\markstar\ 对称性]
\begin{enumerate}[label=(\arabic*)]
    \item 积分区域 $D$ 关于 $x$ 轴对称.
    \begin{enumerate}[label=(\roman*)]
        \item 若 $f(x,y)=-f(x,-y)$, 则 $\displaystyle\iint\limits_Df(x,y)\,\mathrm{d}x\,\mathrm{d}y=0$;
        \item 若 $f(x,y)=f(x,-y)$, 则
        \[
            \iint\limits_Df(x,y)\,\mathrm{d}x\,\mathrm{d}y
            =2\iint\limits_{D\cap\{y\geq0\}}f(x,y)\,\mathrm{d}x\,\mathrm{d}y.
        \]
    \end{enumerate}

    \item 积分区域 $D$ 关于 $y$ 轴对称.
    \begin{enumerate}[label=(\roman*)]
        \item 若 $f(x,y)=-f(-x,y)$, 则 $\displaystyle\iint\limits_Df(x,y)\,\mathrm{d}x\,\mathrm{d}y=0$;
        \item 若 $f(x,y)=f(-x,y)$, 则
        \[
            \iint\limits_Df(x,y)\,\mathrm{d}x\,\mathrm{d}y
            =2\iint\limits_{D\cap\{x\geq0\}}f(x,y)\,\mathrm{d}x\,\mathrm{d}y.
        \]
    \end{enumerate}

    \item 积分区域 $D$ 关于原点对称.
    \begin{enumerate}[label=(\roman*)]
        \item 若 $f(x,y)=-f(-x,-y)$, 则 $\displaystyle\iint\limits_Df(x,y)\,\mathrm{d}x\,\mathrm{d}y=0$;
        \item 若 $f(x,y)=f(-x,-y)$, 则
        \[
            \iint\limits_Df(x,y)\,\mathrm{d}x\,\mathrm{d}y
            =2\iint\limits_{D_1}f(x,y)\,\mathrm{d}x\,\mathrm{d}y,
        \]
        其中 $D_1$ 为区域 $D$ 的一半.
    \end{enumerate}
\end{enumerate}
\end{theorem}

\begin{remark}
对称性一般是用来简化计算, 消去部分被积函数项, 缩小积分区域.
\end{remark}

\begin{example}
计算积分
\[
    I=\iint\limits_Dx\bigl(1+yf(x^2+y^2)\bigr)\,\mathrm{d}x\,\mathrm{d}y,
\]
其中 $D$ 是由 $y=x^3,y=1,x=-1$ 所围成的区域, $f(x)$ 是实值连续函数.
\end{example}
\exampleblank

\begin{example}
计算积分:
\begin{enumerate}[label=(\arabic*)]
    \item
    \[
        \iint\limits_D\left|xy-\dfrac{1}{4}\right|\,\mathrm{d}x\,\mathrm{d}y,
        \qquad D=\lbrack0,1\rbrack\times\lbrack0,1\rbrack;
    \]
    \item
    \[
        \iint\limits_{|x|\leq1,\ 0\leq y\leq2}\sqrt{|y-x^2|}\,\mathrm{d}x\,\mathrm{d}y.
    \]
\end{enumerate}
\end{example}
\exampleblank

\begin{example}
计算积分:
\begin{enumerate}[label=(\arabic*)]
    \item $\displaystyle\iint\limits_{x^2+y^2\leq5}\operatorname{sgn}(x^2-y^2+3)\,\mathrm{d}x\,\mathrm{d}y$;
    \item $\displaystyle\iint\limits_{0\leq x,y\leq2}[x+y]\,\mathrm{d}x\,\mathrm{d}y$.
\end{enumerate}
\end{example}
\exampleblank

\begin{example}
举例说明下列命题:
\begin{enumerate}[label=(\arabic*)]
    \item 设 $D=\{(x,y)\mid a\leq x\leq b,\ c<y\leq d\}$, $f\in C(D)$, 此时累次积分不一定能交换次序;
    \item 存在 $\lbrack0,1\rbrack\times\lbrack0,1\rbrack$ 上非负不可积函数 $f(x,y)$, 使得
    \[
        \int_0^1\mathrm{d}x\int_0^1f(x,y)\,\mathrm{d}y
        =0
        =\int_0^1\mathrm{d}y\int_0^1f(x,y)\,\mathrm{d}x.
    \]
\end{enumerate}
\end{example}
\exampleblank

\newpage

\subsection{变量替换}

\begin{theorem}[\markstar]
设 $D,\widetilde D$ 是 $\mathbb{R}^2$ 中由逐段光滑曲线围成的有界闭区域, 变换
\[
    T:\begin{cases}
        x=x(u,v),\\
        y=y(u,v)
    \end{cases}
    :\widetilde D\to D
\]
是同胚变换(双射及 $T,T^{-1}$ 连续), 且 $T\in C^{(2)}(\widetilde D)$, 记
\[
    J=\dfrac{\partial(x,y)}{\partial(u,v)}.
\]
若 $f(x,y)$ 在 $D$ 上可积, 则
\[
    \iint\limits_Df(x,y)\,\mathrm{d}x\,\mathrm{d}y
    =\iint\limits_{\widetilde D}f\bigl(x(u,v),y(u,v)\bigr)|J(u,v)|\,\mathrm{d}u\,\mathrm{d}v.
\]
\end{theorem}

\begin{remark}
特别地,
\[
    \begin{cases}
        x=r\cos\theta,\\
        y=r\sin\theta
    \end{cases}
\]
为极坐标变换, $\theta$ 为沿 $x$ 正半轴逆时针方向的旋转角, $r$ 为到原点的距离, 此时 $J=r$. 换元的目的旨在化简积分区域(化为 $x$ 型或 $y$ 型或二者多个的并)和被积函数.
\end{remark}

\begin{example}
计算积分:
\begin{enumerate}[label=(\arabic*)]
    \item
    \[
        I=\iint\limits_D\dfrac{3x}{y^2+xy}\,\mathrm{d}x\,\mathrm{d}y,
    \]
    其中 $D$ 为平面曲线 $xy=1,xy=3,y^2=x,y^2=3x$ 所围成的区域;
    \item
    \[
        J=\iint\limits_D\dfrac{(\sqrt{x}+\sqrt{y})^4}{x^2}\,\mathrm{d}x\,\mathrm{d}y,
    \]
    其中 $D$ 为由 $x$ 轴、$y=x$、$\sqrt{x}+\sqrt{y}=1$ 和 $\sqrt{x}+\sqrt{y}=2$ 所围成的区域.
\end{enumerate}
\end{example}
\exampleblank

\begin{example}
计算积分:
\begin{enumerate}[label=(\arabic*)]
    \item
    \[
        \iint\limits_D\dfrac{x^2-y^2}{\sqrt{x+y+3}}\,\mathrm{d}x\,\mathrm{d}y,
        \qquad D=\{(x,y)\mid |x|+|y|\leq1\};
    \]
    \item
    \[
        \iint\limits_D\dfrac{(x+y)\ln(1+xy)}{\sqrt{1-x-y}}\,\mathrm{d}x\,\mathrm{d}y,
        \qquad D=\left\{(x,y)\mid0\leq y\leq x,\ \dfrac{3}{4}\leq x+y\leq1\right\};
    \]
    \item
    \[
        \iint\limits_De^{\frac{x-y}{x+y}}\,\mathrm{d}x\,\mathrm{d}y,
    \]
    其中 $D$ 是 $x=0,y=0,x+y=1$ 所围成的有界闭区域.
\end{enumerate}
\end{example}
\exampleblank

\begin{example}
计算积分:
\begin{enumerate}[label=(\arabic*)]
    \item $\displaystyle\iint\limits_D(x+y)\,\mathrm{d}x\,\mathrm{d}y$, $D$ 是由曲线 $x^2+y^2=x+y$ 所围成的区域;
    \item
    \[
        \iint\limits_D\dfrac{1}{xy}\,\mathrm{d}x\,\mathrm{d}y,
    \]
    其中 $D$ 为平面上由
    \[
        2\leq\dfrac{x}{x^2+y^2},\qquad
        \dfrac{y}{x^2+y^2}\leq4
    \]
    确定的区域;
    \item
    \[
        \iint\limits_D\bigl(3x^3+x^2+y^2+2x-2y+1\bigr)\,\mathrm{d}x\,\mathrm{d}y,
    \]
    其中
    \[
        D=\{(x,y)\mid1\leq x^2+(y-1)^2\leq2,\ x^2+y^2\leq1\}.
    \]
\end{enumerate}
\end{example}
\exampleblank

\begin{example}
计算积分:
\begin{enumerate}[label=(\arabic*)]
    \item
    \[
        \iint\limits_{x^2+y^2\leq\frac{3}{16}}
        \min\left\{\sqrt[3]{\dfrac{3}{16}-x^2-y^2},\ 2(x^2+y^2)\right\}\,\mathrm{d}x\,\mathrm{d}y;
    \]
    \item
    \[
        \iint\limits_{x^2+y^2\leq1}\left|\dfrac{x+y}{\sqrt{2}}-x^2-y^2\right|\,\mathrm{d}x\,\mathrm{d}y.
    \]
\end{enumerate}
\end{example}
\exampleblank

\begin{example}
计算积分:
\begin{enumerate}[label=(\arabic*)]
    \item $\displaystyle\iint\limits_{x^2+4y^2\leq1}(x^2+y^2)\,\mathrm{d}x\,\mathrm{d}y$;
    \item $\displaystyle\iint\limits_{\sqrt{x}+\sqrt{y}\leq1,\ x,y\geq0}xy\,\mathrm{d}x\,\mathrm{d}y$;
    \item $\displaystyle\iint\limits_{x^4+y^4\leq1}(x^2+y^2)\,\mathrm{d}x\,\mathrm{d}y$.
\end{enumerate}
\end{example}
\exampleblank

\begin{example}
计算积分
\[
    \iint\limits_D\sqrt{\dfrac{xy}{x+y}}\,\mathrm{d}x\,\mathrm{d}y,
\]
其中 $D$ 是由 $x=0,y=0,x+y=1$ 所围成的区域.
\end{example}
\exampleblank

\newpage

\subsection{理论证明}

由一元函数积分中值定理证明过程易知, 重积分中值定理仍成立.

\begin{theorem}[\markstar]
设 $f(x,y)$ 在有界闭区域 $D$ 上连续, 则存在 $(\xi,\eta)\in D$, 使得
\[
    \iint\limits_Df(x,y)\,\mathrm{d}x\,\mathrm{d}y=f(\xi,\eta)|D|.
\]
其中 $|D|$ 表示 $D$ 的面积.
\end{theorem}

\begin{remark}
事实上可以证明上述中值点 $(\xi,\eta)$ 可以有无穷多个.
\end{remark}

\begin{example}
设 $f(x,y)\in C^{(2)}(D)$, $D=\lbrack a,b\rbrack\times\lbrack c,d\rbrack$, 则
\begin{enumerate}[label=(\arabic*)]
    \item
    \[
        \iint\limits_D f_{xy}(x,y)\,\mathrm{d}x\,\mathrm{d}y
        =\iint\limits_D f_{yx}(x,y)\,\mathrm{d}x\,\mathrm{d}y;
    \]
    \item 利用 (1) 证明 $f_{xy}(x,y)=f_{yx}(x,y)$.
\end{enumerate}
\end{example}
\exampleblank

\begin{example}
设 $f\in C(\mathbb{R})$, $h>0$,
\[
    F(x)=\dfrac{1}{h^2}\int_0^h\mathrm{d}y\int_0^h f(x+y+z)\,\mathrm{d}z,
\]
求 $F''(x)$.
\end{example}
\exampleblank

\begin{example}
设 $x=x(u,v),y=y(u,v)$ 有连续偏导数, 一一对应地将区域 $D'$ 映射到 $xOy$ 平面的区域 $D$, 满足 $J\ne0$, 且
\[
    \dfrac{\partial x}{\partial u}=\dfrac{\partial y}{\partial v},\qquad
    \dfrac{\partial x}{\partial v}=-\dfrac{\partial y}{\partial u},
\]
试证
\[
    \iint\limits_D\left[\left(\dfrac{\partial f}{\partial x}\right)^2+
    \left(\dfrac{\partial f}{\partial y}\right)^2\right]\mathrm{d}x\,\mathrm{d}y
    =\iint\limits_{D'}\left[\left(\dfrac{\partial f}{\partial u}\right)^2+
    \left(\dfrac{\partial f}{\partial v}\right)^2\right]\mathrm{d}u\,\mathrm{d}v.
\]
\end{example}
\exampleblank

\begin{example}
设 $f(x,y)\geq0$ 在 $D:x^2+y^2\leq a^2$ 上有连续的一阶偏导数, 边界上取值为零, 证明
\[
    \left|\iint\limits_Df(x,y)\,\mathrm{d}x\,\mathrm{d}y\right|
    \leq\dfrac{\pi a^3}{3}
    \max_{(x,y)\in D}\sqrt{\left(\dfrac{\partial f}{\partial x}\right)^2+
    \left(\dfrac{\partial f}{\partial y}\right)^2}.
\]
\end{example}
\exampleblank

\begin{example}
设 $D=\{(x,y)\mid0\leq x,y\leq1\}$, $f(x,y)$ 在 $D$ 上四次连续可微, 当 $(x,y)\in\partial D$ 时 $f(x,y)=0$. 对任意 $(x,y)\in D$, 有
\[
    \left|\dfrac{\partial^4f(x,y)}{\partial x^2\partial y^2}\right|\leq M,
\]
则
\[
    \left|\iint\limits_Df(x,y)\,\mathrm{d}x\,\mathrm{d}y\right|\leq\dfrac{M}{144}.
\]
\end{example}
\exampleblank

\begin{example}[\markstar]
利用高维积分证明低维积分不等式:
\begin{enumerate}[label=(\arabic*)]
    \item (Minkowski 不等式)设 $f(x,y)$ 在 $\mathbb{R}^2$ 上连续, 则
    \[
        \left\{\int_a^b\left[\int_c^d f(x,y)\,\mathrm{d}y\right]^2\mathrm{d}x\right\}^{\frac{1}{2}}
        \leq
        \int_c^d\left\{\int_a^b f^2(x,y)\,\mathrm{d}x\right\}^{\frac{1}{2}}\mathrm{d}y.
    \]
    \item 设 $f(x)$ 是 $\lbrack0,1\rbrack$ 上递减正值函数, 则
    \[
        \dfrac{\displaystyle\int_0^1x f^2(x)\,\mathrm{d}x}
        {\displaystyle\int_0^1x f(x)\,\mathrm{d}x}
        \leq
        \dfrac{\displaystyle\int_0^1 f^2(x)\,\mathrm{d}x}
        {\displaystyle\int_0^1 f(x)\,\mathrm{d}x}.
    \]
    \item 设 $p(x)\in R\lbrack a,b\rbrack$, $p(x)\geq0$, $f(x),g(x)$ 在 $\lbrack a,b\rbrack$ 上递增, 则
    \[
        \left(\int_a^bp(x)f(x)\,\mathrm{d}x\right)
        \left(\int_a^bp(x)g(x)\,\mathrm{d}x\right)
        \leq
        \left(\int_a^bp(x)\,\mathrm{d}x\right)
        \left(\int_a^bp(x)f(x)g(x)\,\mathrm{d}x\right).
    \]
\end{enumerate}
\end{example}
\exampleblank

\newpage

\subsection{二重积分的应用}

\begin{example}
求由下列曲面围成的立体体积 $V$:
\begin{enumerate}[label=(\arabic*)]
    \item $z=0,z=x^2+y^2,x^2+y^2=x,x^2+y^2=2x$;
    \item $z=x^2+2y^2,z=2-x^2$;
    \item $\dfrac{x^2}{a^2}+\dfrac{y^2}{b^2}+\dfrac{z^2}{c^2}=1$, $\dfrac{x^2}{a^2}+\dfrac{y^2}{b^2}=\dfrac{z^2}{c^2}$, $z>0$, $a,b,c>0$.
\end{enumerate}
\end{example}
\exampleblank

\begin{example}
求下列曲线围成的区域面积 $S$:
\begin{enumerate}[label=(\arabic*)]
    \item $x=0,y=0,\sqrt[4]{\dfrac{x}{a}}+\sqrt[4]{\dfrac{y}{b}}=1$, $a,b>0$;
    \item $\left(\dfrac{x^2}{a^2}+\dfrac{y^2}{b^2}\right)^2=\dfrac{xy}{c^2}$, $a,b,c>0$;
    \item $\dfrac{x^2}{a^2}+\dfrac{y^2}{b^2}=hx+ky$, $a,b,h,k>0$.
\end{enumerate}
\end{example}
\exampleblank

\begin{example}[\markstar]
用二重积分计算
\[
    \int_0^{+\infty}e^{-x^2}\,\mathrm{d}x.
\]
\end{example}
\exampleblank

\newpage

\section{三重积分}

本节整体内容与上节差不多, 分为重积分化累次积分、变量替换、理论证明以及应用, 但大部分并没有什么新的东西, 只不过升了一维导致计算更加复杂了些. 值得一提的是, 本节变量替换中多了正交变换与柱坐标变换.

\newpage

\subsection{化为累次积分}

不同于二重积分的换序, 三重积分的换序涉及到立体图形, 可能会因图像复杂而难以刻画积分区域, 进而难以换序. 事实上, 可以换个思路, 不必 $x,y,z$ 三个变量一起换, 可以视 $z$ 为常数, 然后先只换 $x,y$, 此即大大降低了难度, 进而可通过多换几次二元序达到换三元序的目的.

\begin{example}[\markstar]
试将下列三重积分次序 $z\to y\to x$ 改变为 $x\to y\to z$:
\begin{enumerate}[label=(\arabic*)]
    \item $\displaystyle\int_0^1\mathrm{d}x\int_0^{1-x}\mathrm{d}y\int_0^{x+y}f(x,y,z)\,\mathrm{d}z$;
    \item $\displaystyle\int_{-1}^1\mathrm{d}x\int_{-\sqrt{1-x^2}}^{\sqrt{1-x^2}}\mathrm{d}y\int_{\sqrt{x^2+y^2}}^1f(x,y,z)\,\mathrm{d}z$;
    \item $\displaystyle\int_0^1\mathrm{d}x\int_0^x\mathrm{d}y\int_0^{xy}f(x,y,z)\,\mathrm{d}z$.
\end{enumerate}
\end{example}
\exampleblank

\begin{theorem}[\markstar\ 对称性]
\begin{enumerate}[label=(\arabic*)]
    \item 积分区域 $V$ 关于 $xy$ 平面对称, 则考虑 $f(x,y,z)$ 关于 $z$ 的奇偶性;
    \item 积分区域 $V$ 关于 $yz$ 平面对称, 则考虑 $f(x,y,z)$ 关于 $x$ 的奇偶性;
    \item 积分区域 $V$ 关于 $xz$ 平面对称, 则考虑 $f(x,y,z)$ 关于 $y$ 的奇偶性;
    \item 积分区域 $V$ 关于原点对称, 则考虑 $f(x,y,z)$ 关于原点的对称性.
\end{enumerate}
相应地, 奇对称项在对称区域上的三重积分为 $0$, 偶对称项的积分可化为半区域积分的 $2$ 倍.
\end{theorem}

\begin{example}
计算积分:
\begin{enumerate}[label=(\arabic*)]
    \item
    \[
        \iiint\limits_V\dfrac{z\ln(x^2+y^2+z^2+1)}{x^2+y^2+z^2+1}\,\mathrm{d}x\,\mathrm{d}y\,\mathrm{d}z,
    \]
    其中 $V:x^2+y^2+z^2\leq1$;
    \item
    \[
        \iiint\limits_{\substack{x^2+y^2+z^2\leq1\\x^2+y^2-z^2\geq\frac{1}{2}}}z\,\mathrm{d}x\,\mathrm{d}y\,\mathrm{d}z.
    \]
\end{enumerate}
\end{example}
\exampleblank

\begin{example}
求积分
\[
    \iiint\limits_V(x+y+z)\,\mathrm{d}x\,\mathrm{d}y\,\mathrm{d}z,
\]
其中 $V$ 是由平面 $x+y+z=1$ 以及三个坐标平面所围成的区域.
\end{example}
\exampleblank

\begin{example}
求区域
\[
    V:\quad0\leq x\leq1,\quad0\leq y\leq x,\quad x+y\leq z\leq e^{x+y}
\]
的体积.
\end{example}
\exampleblank

\begin{example}
设 $\Omega$ 由 $z^2=x^2+y^2,z=0,xy=1,xy=2,y=3x,y=4x$ 所围成, 求积分
\[
    \iiint\limits_\Omega x^2y^2z\,\mathrm{d}x\,\mathrm{d}y\,\mathrm{d}z.
\]
\end{example}
\exampleblank

\begin{example}
计算积分
\[
    \iiint\limits_V\dfrac{\mathrm{d}V}{y^2+z^2},
\]
其中 $V$ 为一棱台, 其六个顶点为 $A(0,0,1)$, $B(0,1,1)$, $C(1,1,1)$, $D(0,0,2)$, $E(0,2,2)$, $F(2,2,2)$.
\end{example}
\exampleblank

\begin{example}
计算
\[
    \iiint\limits_V y\cos(x+z)\,\mathrm{d}x\,\mathrm{d}y\,\mathrm{d}z,
\]
其中 $V$ 是由 $y=\sqrt{x},x+z=\dfrac{\pi}{2},y=0,z=0$ 所围成的区域.
\end{example}
\exampleblank

\begin{example}
计算
\[
    \iiint\limits_V(xy+yz+zx)\,\mathrm{d}x\,\mathrm{d}y\,\mathrm{d}z,
\]
其中 $V:x\geq0,y\geq0,x^2+y^2\leq1,0\leq z\leq1$.
\end{example}
\exampleblank

\newpage

\subsection{变量替换}

\begin{theorem}[\markstar]
设 $\widetilde\Omega,\Omega$ 是 $(x,y,z)$ 空间中的有界闭区域, 变换
\[
    T:\begin{cases}
        x=x(u,v,w),\\
        y=y(u,v,w),\\
        z=z(u,v,w)
    \end{cases}
    :\widetilde\Omega\to\Omega
\]
是同胚变换, 且 $T\in C^{(2)}(\widetilde\Omega)$,
\[
    J=\dfrac{\partial(x,y,z)}{\partial(u,v,w)}.
\]
若 $f(x,y,z)$ 在 $\Omega$ 上可积, 则
\[
    \iiint\limits_\Omega f(x,y,z)\,\mathrm{d}x\,\mathrm{d}y\,\mathrm{d}z
    =\iiint\limits_{\widetilde\Omega}
    f\bigl(x(u,v,w),y(u,v,w),z(u,v,w)\bigr)|J|\,\mathrm{d}u\,\mathrm{d}v\,\mathrm{d}w.
\]
\end{theorem}

\begin{remark}
常用特例包括:
\begin{enumerate}[label=(\arabic*)]
    \item 正交变换
    \[
        \begin{cases}
            x=a_{11}u+a_{12}v+a_{13}w,\\
            y=a_{21}u+a_{22}v+a_{23}w,\\
            z=a_{31}u+a_{32}v+a_{33}w,
        \end{cases}
    \]
    其中 $A=(a_{ij})$ 为正交矩阵, 此时 $|J|=1$;
    \item 柱坐标变换
    \[
        x=r\cos\theta,\qquad y=r\sin\theta,\qquad z=z,
    \]
    其中 $0\leq r<+\infty,0\leq\theta<2\pi,-\infty<z<+\infty$, 此时 $J=r$;
    \item 球坐标变换
    \[
        x=r\sin\varphi\cos\theta,\qquad
        y=r\sin\varphi\sin\theta,\qquad
        z=r\cos\varphi,
    \]
    其中 $0\leq r<+\infty,0\leq\varphi\leq\pi,0\leq\theta<2\pi$, 此时 $J=r^2\sin\varphi$.
\end{enumerate}
\end{remark}

\begin{example}
\begin{enumerate}[label=(\arabic*)]
    \item 证明
    \[
        \iint\limits_S f(ax+by)\,\mathrm{d}x\,\mathrm{d}y
        =2\int_{-1}^1\sqrt{1-u^2}\,f\left(u\sqrt{a^2+b^2}\right)\mathrm{d}u,
    \]
    其中 $S:x^2+y^2\leq1,a^2+b^2\ne0$;
    \item 证明
    \[
        \iiint\limits_Vf(ax+by+cz)\,\mathrm{d}x\,\mathrm{d}y\,\mathrm{d}z
        =\pi\int_{-1}^1(1-u^2)f\left(u\sqrt{a^2+b^2+c^2}\right)\mathrm{d}u,
    \]
    其中 $V:x^2+y^2+z^2\leq1,a^2+b^2+c^2\ne0$.
\end{enumerate}
\end{example}
\exampleblank

\begin{example}
计算积分
\[
    \iiint\limits_\Omega x^2\sqrt{x^2+y^2}\,\mathrm{d}x\,\mathrm{d}y\,\mathrm{d}z,
\]
其中 $\Omega$ 是曲面 $z=\sqrt{x^2+y^2}$ 与 $z=x^2+y^2$ 围成的有界闭区域.
\end{example}
\exampleblank

\begin{example}
\begin{enumerate}[label=(\arabic*)]
    \item 计算椭圆
    \[
        (a_1x+b_1y+c_1)^2+(a_2x+b_2y+c_2)^2=1,
        \qquad a_1b_2-b_1a_2\ne0
    \]
    的面积;
    \item 计算椭球
    \[
        (a_1x+b_1y+c_1z)^2+(a_2x+b_2y+c_2z)^2+(a_3x+b_3y+c_3z)^2=1
    \]
    的体积;
    \item 计算由平面
    \[
        a_1x+b_1y+c_1z=\pm1,\quad
        a_2x+b_2y+c_2z=\pm1,\quad
        a_3x+b_3y+c_3z=\pm1
    \]
    所围立体的体积.
\end{enumerate}
\end{example}
\exampleblank

\begin{example}
计算积分
\[
    \iiint\limits_\Omega z^2\,\mathrm{d}x\,\mathrm{d}y\,\mathrm{d}z,
\]
其中 $\Omega$ 是 $x^2+y^2+z^2\leq a^2$ 与 $x^2+y^2+(z-a)^2\leq a^2$ $(a>0)$ 的公共部分.
\end{example}
\exampleblank

\begin{example}
计算积分
\[
    \iiint\limits_{x^2+y^2+z^2\leq1}\cos(ax+by+cz)\,\mathrm{d}x\,\mathrm{d}y\,\mathrm{d}z.
\]
\end{example}
\exampleblank

\begin{example}
计算积分:
\begin{enumerate}[label=(\arabic*)]
    \item $\displaystyle\iiint\limits_V(x^2+y^2)\,\mathrm{d}V$, 其中 $V$ 是由 $x^2+y^2+(z-2)^2\geq4$, $x^2+y^2+(z-1)^2\leq9$, $z\geq0$ 所围成的空心立体;
    \item
    \[
        \iiint\limits_V\dfrac{xyz}{x^2+y^2}\,\mathrm{d}x\,\mathrm{d}y\,\mathrm{d}z,
    \]
    其中 $V$ 是曲面 $(x^2+y^2+z^2)^2=a^2xy$ $(a>0)$ 位于第一象限所围立体;
    \item
    \[
        \iiint\limits_V(x+y+z)^2\,\mathrm{d}x\,\mathrm{d}y\,\mathrm{d}z,
    \]
    其中 $V$ 为 $x^2+y^2\leq2az$, $x^2+y^2+z^2\leq3$ 所围立体.
\end{enumerate}
\end{example}
\exampleblank

\begin{example}[\markstar]
\begin{enumerate}[label=(\arabic*)]
    \item 设 $a,b,c>0$, 证明
    \[
        \iiint\limits_Vx^{a-1}y^{b-1}z^{c-1}\,\mathrm{d}x\,\mathrm{d}y\,\mathrm{d}z
        =\dfrac{1}{a+b+c}\dfrac{\Gamma(a)\Gamma(b)\Gamma(c)}{\Gamma(a+b+c)},
    \]
    其中 $V$ 为四面体 $x\geq0,y\geq0,z\geq0,x+y+z\leq1$;
    \item 计算
    \[
        \iiint\limits_Vxy^9z^8w^4\,\mathrm{d}x\,\mathrm{d}y\,\mathrm{d}z,
        \qquad w=1-x-y-z,
    \]
    其中 $V:x,y,z\geq0,x+y+z\leq1$.
\end{enumerate}
\end{example}
\exampleblank

\newpage

\subsection{理论及应用}

\begin{example}
设 $f(u)$ 在 $\lbrack0,+\infty)$ 上可微, 试求极限
\[
    \lim\limits_{t\to0^+}\dfrac{1}{\pi t^4}
    \iiint\limits_{x^2+y^2+z^2\leq t^2}
    f\left(\sqrt{x^2+y^2+z^2}\right)\,\mathrm{d}x\,\mathrm{d}y\,\mathrm{d}z.
\]
\end{example}
\exampleblank

\begin{example}[\markstar]
证明
\[
    \lim\limits_{n\to\infty}\dfrac{1}{n^4}
    \iiint\limits_{r\leq n}[r]\,\mathrm{d}x\,\mathrm{d}y\,\mathrm{d}z=\pi,
    \qquad r=\sqrt{x^2+y^2+z^2}.
\]
\end{example}
\exampleblank

\begin{example}
设 $f(x,y,z)$ 在
\[
    V=\{(x,y,z)\mid0\leq x,y,z\leq1\}
\]
上有六阶连续偏导数, $f$ 在边界上恒为零, 且
\[
    \left|\dfrac{\partial^6f(x,y,z)}{\partial x^2\partial y^2\partial z^2}\right|\leq M,
\]
证明
\[
    \left|\iiint\limits_Vf(x,y,z)\,\mathrm{d}x\,\mathrm{d}y\,\mathrm{d}z\right|
    \leq\dfrac{1}{8}\cdot\dfrac{M}{6^3}.
\]
\end{example}
\exampleblank

\begin{example}
求由下列曲面围成的立体的体积 $V$:
\begin{enumerate}[label=(\arabic*)]
    \item $\left(\dfrac{x^2}{a^2}+\dfrac{y^2}{b^2}+\dfrac{z^2}{c^2}\right)^2=hx$;
    \item $(x^2+y^2)^2+z^4=y$;
    \item $x^2+z^2=a^2$, $x^2+z^2=b^2$, $x^2-y^2-z^2=0$, $x>0$, $a<b$.
\end{enumerate}
\end{example}
\exampleblank

\begin{example}[\markstar]
设
\[
    \sum\limits_{i,j=1}^3a_{ij}x_ix_j
\]
表示变量 $(x_1,x_2,x_3)$ 的二次型, 其中 $A=(a_{ij})$ 正定. 证明椭球面
\[
    S:\quad\sum\limits_{i,j=1}^3a_{ij}x_ix_j=1
\]
所包围的体积等于
\[
    \dfrac{4}{3}\pi(\det A)^{-\frac{1}{2}}.
\]
\end{example}
\exampleblank

\newpage

\subsection{\texorpdfstring{$n$}{n} 重积分 *}

$n$ 重积分处理技巧基本上类似于二、三重积分, 由于 $n$ 维的一般性, 计算时多了一种递推方法.

\begin{theorem}[\markstar\ 化为累次积分]
\begin{enumerate}[label=(\arabic*)]
    \item 设
    \[
        V=\{(x_1,x_2,\ldots,x_n)\mid a_i\leq x_i\leq b_i,\ i=1,2,\ldots,n\},
    \]
    $f$ 在 $V$ 上连续, 则
    \[
        \int\cdots\int\limits_V f(x_1,\ldots,x_n)\,\mathrm{d}x_1\cdots\mathrm{d}x_n
        =\int_{a_1}^{b_1}\mathrm{d}x_1\cdots\int_{a_n}^{b_n}f(x_1,\ldots,x_n)\,\mathrm{d}x_n.
    \]

    \item 设
    \[
    \begin{aligned}
        V=\{(x_1,\ldots,x_n)\mid{}&a_1\leq x_1\leq b_1,\ a_2(x_1)\leq x_2\leq b_2(x_1),\ldots,\\
        &a_n(x_1,\ldots,x_{n-1})\leq x_n\leq b_n(x_1,\ldots,x_{n-1})\},
    \end{aligned}
    \]
    $f$ 在 $V$ 上连续, 则相应的 $n$ 重积分可按上述边界依次化为累次积分.
\end{enumerate}
\end{theorem}

换元结论类似. 值得一提的是推广的球坐标变换:
\[
\begin{cases}
    x_1=r\cos\varphi_1,\\
    x_2=r\sin\varphi_1\cos\varphi_2,\\
    x_3=r\sin\varphi_1\sin\varphi_2\cos\varphi_3,\\
    \cdots,\\
    x_{n-1}=r\sin\varphi_1\cdots\sin\varphi_{n-2}\cos\varphi_{n-1},\\
    x_n=r\sin\varphi_1\cdots\sin\varphi_{n-2}\sin\varphi_{n-1}.
\end{cases}
\]
其中 $0\leq r<+\infty$, $0\leq\varphi_1,\ldots,\varphi_{n-2}\leq\pi$, $0\leq\varphi_{n-1}<2\pi$, Jacobi 行列式为
\[
    J=r^{n-1}\sin^{n-2}\varphi_1\sin^{n-3}\varphi_2\cdots\sin^2\varphi_{n-3}\sin\varphi_{n-2}.
\]

\begin{example}
设 $a_1,\ldots,a_n>0$, 又
\[
    S_n(a_1,\ldots,a_n)=\left\{(x_1,\ldots,x_n)\ \middle|\ \dfrac{|x_1|}{a_1}+\cdots+\dfrac{|x_n|}{a_n}\leq1\right\},
\]
计算 $S_n(a_1,\ldots,a_n)$ 的体积.
\end{example}
\exampleblank

\begin{example}
求
\[
    \iiiint\limits_V\sqrt{\dfrac{1-x^2-y^2-z^2-u^2}{1+x^2+y^2+z^2+u^2}}\,
    \mathrm{d}x\,\mathrm{d}y\,\mathrm{d}z\,\mathrm{d}u,
\]
其中 $V:x,y,z,u\geq0$, $x^2+y^2+z^2+u^2\leq1$.
\end{example}
\exampleblank

\begin{example}
计算
\[
    \iiiint\limits_De^{-\langle A\bm{x},\bm{x}\rangle}\,
    \mathrm{d}x_1\,\mathrm{d}x_2\,\mathrm{d}x_3\,\mathrm{d}x_4,
\]
其中
\[
    \langle A\bm{x},\bm{x}\rangle=\sum\limits_{i,j=1}^4a_{ij}x_ix_j
\]
是正定二次型, $D$ 是区域 $\langle A\bm{x},\bm{x}\rangle\leq1$.
\end{example}
\exampleblank

\begin{example}
试求 $\mathbb{R}^n$ 中下列立体 $\Omega$ 的体积 $V$:
\begin{enumerate}[label=(\arabic*)]
    \item
    \[
        \Omega_n(a)=\{(x_1,\ldots,x_n)\mid x_i\geq0,\ x_1+\cdots+x_n\leq a\};
    \]
    \item
    \[
        \Omega=\left\{(x_1,\ldots,x_n)\ \middle|\ x_1^{a_1}+\cdots+x_n^{a_n}\leq1,\ x_i\geq0\right\},
    \]
    其中 $a_i>0$;
    \item
    \[
        \Omega=\{(x_1,\ldots,x_n)\mid x_1^2+x_2^2+\cdots+x_n^2\leq a^2\},\qquad a>0.
    \]
\end{enumerate}
\end{example}
\exampleblank

\begin{example}
设 $f(x)$ 在 $\lbrack a,b\rbrack$ 上连续, 证明: 对任意 $x\in(a,b)$,
\[
    \int_a^x\mathrm{d}x_1\int_a^{x_1}\mathrm{d}x_2\cdots\int_a^{x_n}f(x_{n+1})\,\mathrm{d}x_{n+1}
    =\dfrac{1}{n!}\int_a^x(x-y)^nf(y)\,\mathrm{d}y.
\]
\end{example}
\exampleblank

\begin{example}
设 $f\in C\lbrack0,1\rbrack$, 且
\[
    \int_0^1f(x)\,\mathrm{d}x=A,
\]
求积分
\[
    \int_0^1\cdots\int_0^1
    \dfrac{x_1+x_2+\cdots+x_k}{x_1+x_2+\cdots+x_n}
    f(x_1)f(x_2)\cdots f(x_n)
    \,\mathrm{d}x_1\cdots\mathrm{d}x_n,
    \qquad k\leq n.
\]
\end{example}
\exampleblank

\newpage

\section{反常重积分 *}

\begin{definition}[无界区域]
设 $D$ 是 $\mathbb{R}^2$ 中的无界区域, 其边界由有限条光滑或逐段光滑曲线组成. 函数 $f$ 定义在 $D$ 上, 且在 $D$ 内的任何可求面积的有界子区域上可积. 设 $D_r$ 是 $D$ 内的任一可求面积的有界子区域, 包含 $D\cap B_r$, 其中 $B_r$ 是 $\mathbb{R}^2$ 中以 $r$ 为半径的闭圆盘. 若极限
\[
    I=\lim\limits_{r\to+\infty}\iint\limits_{D_r}f(x,y)\,\mathrm{d}x\,\mathrm{d}y
\]
存在且有限, 并与 $D_r$ 的取法无关, 则称 $f$ 在 $D$ 上积分收敛, 否则称积分发散.
\end{definition}

\begin{remark}
因此, 证明发散的话, 可以找两种收敛于不同值的分割方法, 或一种极限不存在的分割方式.
\end{remark}

\begin{theorem}[\markstar]
设 $f(x,y)$ 是无界区域 $D\subseteq\mathbb{R}^2$ 上的非负函数, 且在 $D$ 内的任何可求面积的有界子区域上可积, $D_n$ 是包含 $D\cap B_n$ 的一列可求面积的有界子区域, 则
\[
    \iint\limits_D f(x,y)\,\mathrm{d}x\,\mathrm{d}y\text{ 收敛}
    \iff
    \left\{\iint\limits_{D_n}f(x,y)\,\mathrm{d}x\,\mathrm{d}y\right\}\text{ 有界}
\]
且等价于
\[
    \lim\limits_{n\to+\infty}\iint\limits_{D_n}f(x,y)\,\mathrm{d}x\,\mathrm{d}y
\]
存在.
\end{theorem}

\begin{remark}
即非负函数时, 只需找到一种特殊的区间套, 使得 $f(x,y)$ 在其上积分有界即可.
\end{remark}

\begin{definition}[无界函数]
设 $D$ 为 $\mathbb{R}^2$ 上可求面积的有界区域, 点 $P_0\in D$, 函数 $f$ 定义在 $D\setminus\{P_0\}$ 上且对任何点 $P_0$ 的可求面积的邻域 $\Delta$, $f$ 在 $D\setminus\Delta$ 上有界可积. 如果
\[
    I=\lim\limits_{d(\Delta)\to0}\iint\limits_{D\setminus\Delta}f(x,y)\,\mathrm{d}x\,\mathrm{d}y
\]
存在且有界, 并与 $\Delta$ 的取法无关, 其中 $d(\Delta)$ 为 $\Delta$ 的直径, 则称 $f$ 在 $D$ 上积分收敛, 否则称积分发散.
\end{definition}

\begin{remark}
显然, 无界函数也有相应于前述定理的非负函数充要条件版本. 类似于一元版本, 对于非负函数有相应比较判别法. 特别地, 当取 $g(x,y)=\dfrac{C}{r^p}$, $r=\sqrt{x^2+y^2}$ 时, 可得到下面的 Cauchy 判别法.
\end{remark}

\begin{theorem}[\markstar\ Cauchy 判别法]
设 $D$ 为无界区域, 其边界由有限条光滑或逐段连续曲线组成, $f$ 在 $D$ 内的任一可求面积的有界子区域上可积, 则:
\begin{enumerate}[label=(\arabic*)]
    \item 若对充分大的 $r$, 有 $|f(x,y)|\leq\dfrac{C}{r^p}$, $p>2$, 则反常积分 $\displaystyle\iint\limits_Df(x,y)\,\mathrm{d}x\,\mathrm{d}y$ 收敛;
    \item 若对充分大的 $r$, 有 $|f(x,y)|>\dfrac{C}{r^p}$, $p\leq2$, 则反常积分 $\displaystyle\iint\limits_Df(x,y)\,\mathrm{d}x\,\mathrm{d}y$ 发散.
\end{enumerate}
\end{theorem}

总的来说, 反常重积分与反常积分大部分理论完全相似, 但是却有如下惊人的差别.

\begin{theorem}[\markstar]
$f$ 反常可积当且仅当 $|f|$ 反常可积.
\end{theorem}

\begin{example}
讨论反常重积分
\[
    I=\iint\limits_D\dfrac{\mathrm{d}x\,\mathrm{d}y}{(x+y)^p}
\]
的敛散性, 其中 $D=\{(x,y)\mid0\leq x\leq1,\ x+y\geq1\}$. 收敛时, 求积分的值.
\end{example}
\exampleblank

\begin{example}
设 $D:1\leq x,y<+\infty$,
\[
    f(x,y)=\dfrac{x^2-y^2}{(x^2+y^2)^2},
\]
令
\[
    I=\iint\limits_Df(x,y)\,\mathrm{d}x\,\mathrm{d}y,
\]
\[
    I_1=\int_1^{+\infty}\mathrm{d}x\int_1^{+\infty}f(x,y)\,\mathrm{d}y,
    \qquad
    I_2=\int_1^{+\infty}\mathrm{d}y\int_1^{+\infty}f(x,y)\,\mathrm{d}x.
\]
则 $I$ 发散, $I_1$ 与 $I_2$ 收敛.
\end{example}
\exampleblank

\begin{example}
证明积分
\[
    I=\iint\limits_D\sin(x^2+y^2)\,\mathrm{d}x\,\mathrm{d}y
\]
不存在, 其中 $D:0\leq x,y<+\infty$.
\end{example}
\exampleblank

\begin{example}
计算下列反常积分:
\begin{enumerate}[label=(\arabic*)]
    \item $\displaystyle I=\iint\limits_D\dfrac{\mathrm{d}x\,\mathrm{d}y}{(1+x+y)^3}$, $D:x\geq0,y\geq0$;
    \item $\displaystyle I=\iint\limits_D\dfrac{\mathrm{d}x\,\mathrm{d}y}{(x+y)^p}$, $D:0\leq x\leq1,x+y\geq1$.
\end{enumerate}
\end{example}
\exampleblank

\begin{example}
判别下列积分的敛散性:
\begin{enumerate}[label=(\arabic*)]
    \item
    \[
        I=\iint\limits_D\dfrac{\varphi(x,y)}{(x^2+y^2)^p}\,\mathrm{d}x\,\mathrm{d}y,
    \]
    其中 $D:x^2+y^2\geq1,x\geq0,y\geq0$, $\varphi\in C(D)$ 且 $m\leq|\varphi(x,y)|\leq M$;
    \item
    \[
        I=\iint\limits_D\dfrac{\mathrm{d}x\,\mathrm{d}y}{x^p+x^q},
        \qquad p,q>0,
    \]
    其中 $D:x\geq0,y\geq0,x+y\geq1$;
    \item
    \[
        I=\iint\limits_D\dfrac{\sin x\sin y}{(x+y)^p}\,\mathrm{d}x\,\mathrm{d}y,
    \]
    其中 $D:x+y\geq1$.
\end{enumerate}
\end{example}
\exampleblank

\begin{example}
设 $D:0\leq x\leq1,0\leq y\leq1,1\leq p<+\infty$, 令
\[
    f(x,y)=
    \begin{cases}
        0,&xy=1,\\
        \dfrac{1}{xy-1},&xy\ne1,
    \end{cases}
\]
则
\[
    \iint\limits_D|f(x,y)|^p\,\mathrm{d}x\,\mathrm{d}y\text{ 收敛}
    \iff1\leq p\leq2.
\]
\end{example}
\exampleblank

\begin{example}
讨论积分
\[
    I=\iiint\limits_D\dfrac{\mathrm{d}x\,\mathrm{d}y\,\mathrm{d}z}{|x|^p+|y|^q+|z|^r},
    \qquad p,q,r>0,
\]
的敛散性, 其中 $D:|x|+|y|+|z|\geq1$.
\end{example}
\exampleblank
